\documentclass[aspectratio=169,11pt]{beamer}
\usepackage[utf8]{inputenc}
\usepackage[T1]{fontenc}
\usepackage[spanish]{babel}
\usepackage{lmodern}
\usepackage{booktabs}
\usepackage{tabularx}
\usepackage{hyperref}
\usepackage{enumitem}

\usetheme{Madrid}
\usecolortheme{default}
\setbeamertemplate{navigation symbols}{}
\setbeamertemplate{footline}[frame number]

\title{IA agéntica para la creación de contenido académico}
\subtitle{Sesión 3 --- Libros/ebooks I: arquitectura y outline}
\author{Juliho Castillo}
\institute{CADi 40\,h · Escuela de Ingeniería y Ciencias\\
Tecnológico de Monterrey}
\date{Clave sugerida: \texttt{CADI-EIC-IA-AGENTICA-01}}

\begin{document}

\begin{frame}
\titlepage
\end{frame}

\begin{frame}{Agenda (3\,h) y producto E3}
\begin{tabularx}{\textwidth}{@{}lclX@{}}
\toprule
Bloque & Tiempo & Contenido \\
\midrule
Recap & 10 min & Brief/orquestación $\rightarrow$ estructura \\
Arquitectura STEM & 30 min & Capítulos, RA, notación, ejercicios \\
Demo outline & 20 min & Brief $\rightarrow$ outline con restricciones \\
Taller & 70 min & Arquitectura + outline + símbolos \\
Pares & 25 min & Coherencia pedagógica \\
Cierre & 15 min & TI avance microartefacto \\
\bottomrule
\end{tabularx}
\vspace{0.5em}
\textbf{E3:} arquitectura de obra + outline documentado con RA.
\end{frame}

\begin{frame}{Resultados de aprendizaje (sesión)}
Al finalizar, usted podrá:
\begin{enumerate}
\item Traducir un brief en \textbf{arquitectura} (capítulos, RA, artefactos)
\item Producir un \textbf{outline} de 6--10 nodos con RA en nodos clave
\item Bosquejar un \textbf{mapa de notación/glosario} STEM
\item Justificar el \textbf{alcance} (qué queda fuera) de un microartefacto viable
\end{enumerate}
\end{frame}

\begin{frame}{Del brief a la obra}
\begin{center}
\texttt{Brief $\rightarrow$ Arquitectura $\rightarrow$ Outline $\rightarrow$ Secciones (ses.\ 4)}
\end{center}
\begin{itemize}
\item El brief fija audiencia y restricciones
\item La arquitectura decide capítulos, artefactos y exclusiones
\item El outline es el \textbf{contrato pedagógico}, no un índice decorativo
\end{itemize}
\begin{block}{Integridad}
Política de uso de IA $\rightarrow$ verificar con \textbf{CEDDIE} \emph{(por edición)}.
\end{block}
\end{frame}

\begin{frame}{Arquitectura de obra académica STEM}
Plantilla mental:
\begin{enumerate}
\item Objetivo y públicos EIC
\item Capítulos / unidades
\item Secciones y artefactos (texto, ejercicios, figuras)
\item Evaluaciones / autoevaluación
\item Glosario y notación
\item \textbf{Qué queda fuera} del microartefacto
\end{enumerate}
\end{frame}

\begin{frame}{Learning outcomes y verbos de Bloom}
\begin{itemize}
\item Asigne RA (verbo + contenido + condición) a nodos clave
\item Prefiera verbos observables: explicar, calcular, diseñar, contrastar\ldots
\item Menos nodos bien alineados $>$ muchos nodos genéricos
\item Cada nodo debe servir a un mensaje o habilidad del curso
\end{itemize}
\begin{alertblock}{Evite}
``Introducción a\ldots'' sin verbo de aprendizaje.
\end{alertblock}
\end{frame}

\begin{frame}{Notación, unidades y glosario}
Antes de redactar secciones:
\begin{itemize}
\item Liste símbolos estables (y sinónimos prohibidos)
\item Fije unidades (SI o las del programa)
\item Defina términos que el agente no debe ``creativizar''
\item Marque qué revisará el agente de consistencia (sesión 4)
\end{itemize}
El mapa de notación ahorra dolor mañana.
\end{frame}

\begin{frame}{Alcance viable del microartefacto}
\begin{itemize}
\item CADi = microartefacto, no el libro completo
\item Elija 1 capítulo o 1--2 secciones profundas
\item Declare exclusiones: anexos, benchmarks, datos no verificados
\item Extensión objetivo: páginas / palabras / \# ejercicios
\end{itemize}
\begin{block}{Diseño 40\,h}
Las 16\,h independientes avanzan el outline hacia la primera sección.
\end{block}
\end{frame}

\begin{frame}{Demo --- brief $\rightarrow$ outline}
En vivo (STEM neutro EIC):
\begin{enumerate}
\item Partir del brief (nivel + restricciones)
\item Generar outline 6--10 nodos
\item Rechazar nodos vagos o fuera de alcance
\item \textbf{HITL:} humano aprueba RA y notación
\end{enumerate}
Observe la diferencia entre índice decorativo y contrato pedagógico.
\end{frame}

\begin{frame}{Ejemplo: outline de 8 nodos}
\begin{enumerate}
\item Gancho / motivación del tema
\item Prerrequisitos y notación
\item Concepto central (RA: explicar)
\item Ejemplo guiado
\item Variante / caso EIC
\item Ejercicio alineado (RA: aplicar)
\item Errores frecuentes
\item Cierre y puente a la siguiente unidad
\end{enumerate}
Adapte; no copie ciegamente.
\end{frame}

\begin{frame}{Roles útiles hoy}
\begin{tabularx}{\textwidth}{@{}lX@{}}
\toprule
Rol & Pregunta \\
\midrule
Planificador & ¿Qué nodos y en qué orden? \\
Pedagogo & ¿Qué RA y verbos? \\
Verificador (preview) & ¿Notación coherente? \\
Humano & ¿Qué queda fuera? \\
\bottomrule
\end{tabularx}
\end{frame}

\begin{frame}{Actividad de taller ($\sim$70 min)}
\textbf{Entregable E3:}
\begin{enumerate}
\item \textbf{Parte A --- Arquitectura} (25 min): objetivo, público, capítulos, exclusiones
\item \textbf{Parte B --- Outline} (30 min): 6--10 nodos + RA en nodos clave + edición humana
\item \textbf{Parte C --- Mapa de notación} (15 min): símbolos / unidades / términos
\end{enumerate}
Nombre: \texttt{Apellido\_Sesion3\_arquitectura-outline}
\end{frame}

\begin{frame}{Rúbrica de coherencia pedagógica}
\begin{tabularx}{\textwidth}{@{}lXX@{}}
\toprule
Criterio & Bien & A mejorar \\
\midrule
Alineación EIC & Audiencia y RA explícitos & Genérico \\
HITL & Nodos editados / \texttt{[VERIFICAR]} & Aceptación pasiva \\
Alcance & Exclusiones claras & Obra infinita \\
Notación & Mapa preliminar & Sin glosario \\
\bottomrule
\end{tabularx}
\end{frame}

\begin{frame}{Revisión por pares ($\sim$25 min)}
Preguntas guía:
\begin{itemize}
\item ¿Cada nodo clave tiene RA observable?
\item ¿Hay nodos que sobran o faltan?
\item ¿La notación está lista para consistencia (ses.\ 4)?
\item ¿El alcance es viable en el tiempo del CADi?
\end{itemize}
\end{frame}

\begin{frame}{Takeaways de la sesión}
\begin{enumerate}
\item Outline = \textbf{contrato pedagógico}, no índice decorativo
\item Menos nodos alineados $>$ muchos genéricos
\item Mapa de notación = inversión para la sesión 4
\item Alcance explícito hace viable el microartefacto
\end{enumerate}
\end{frame}

\begin{frame}{Trabajo independiente --- bloque C}
Antes de la sesión 4:
\begin{itemize}
\item Avanzar el outline hacia la \textbf{primera sección} borrador
\item Refinar mapa de símbolos
\item Traer 1--2 nodos priorizados para generar en aula
\end{itemize}
Parte del presupuesto de \textbf{16\,h} independientes del CADi.
\end{frame}

\begin{frame}{Prep sesión 4 --- Libros/ebooks II}
Próxima sesión (3\,h):
\begin{itemize}
\item Generación supervisada de secciones cortas
\item Ejercicios / figuras alineados a RA
\item Agente de consistencia + diff narrado v0$\rightarrow$v1
\end{itemize}
Traiga E3 y los nodos que expandirá.
\end{frame}

\begin{frame}{Gracias}
\begin{center}
{\Large Juliho Castillo}\\[0.5em]
Escuela de Ingeniería y Ciencias\\[1em]
CADi · IA agéntica para la creación de contenido académico\\[1.5em]
{\small \emph{Verificar políticas de integridad y uso de IA con CEDDIE antes de cada edición.}}
\end{center}
\end{frame}

\end{document}
